% ============================================================
%  sections/04_methodology.tex
%  4. Phương pháp  (~1 trang)
% ============================================================

\section{Phương Pháp}
\label{sec:method}

Chúng tôi đánh giá bốn phương pháp theo thứ tự tăng dần về độ phức tạp
và chi phí tính toán: (1) zero-shot với bộ phát hiện tiền huấn luyện,
(2) tinh chỉnh mô hình đa ngôn ngữ có giám sát, (3) phát hiện thống
kê zero-shot dựa trên perplexity, và (4) học chuyển giao đa ngôn ngữ
(leave-one-language-out).

\subsection{Phương Pháp 1: Baseline Zero-Shot (RoBERTa)}

Chúng tôi áp dụng trực tiếp mô hình \texttt{roberta-base-openai-detector}
~\cite{solaiman2019gpt2detector} --- được tiền huấn luyện để phân biệt
văn bản GPT-2 với văn bản con người tiếng Anh --- lên toàn bộ bộ kiểm
tra đa ngôn ngữ mà không có bất kỳ thích nghi nào. Nhãn dự đoán được
xác định bằng ngưỡng $\tau = 0.5$ trên điểm sigmoid của logit:
\[
\hat{y} = \mathbf{1}\!\left[\sigma(z_{\text{machine}} - z_{\text{human}}) \geq 0.5\right]
\]
Đây là kịch bản ``worst-case'' phản ánh thực tế thường gặp khi triển
khai một detector tiếng Anh cho dữ liệu đa ngôn ngữ.

\subsection{Phương Pháp 2: Tinh Chỉnh XLM-RoBERTa}

Chúng tôi tinh chỉnh \texttt{xlm-roberta-base}~\cite{conneau2020xlmr}
(278M tham số, được tiền huấn luyện trên 100 ngôn ngữ) theo bài toán
phân loại nhị phân (human/machine) trên toàn bộ 12.789 mẫu huấn luyện
đa ngôn ngữ. Các siêu tham số chính:

\begin{itemize}
    \item Số epoch: 3; Learning rate: $2 \times 10^{-5}$; Batch size: 16
    \item Max sequence length: 256 tokens; Precision: fp16
    \item Fallback OOM: batch size 8 + gradient accumulation 2
    \item Metric tối ưu: Macro-F1
\end{itemize}

Mô hình được huấn luyện trên Google Colab T4 (16 GB VRAM). Tổng thời
gian huấn luyện: \textbf{7 phút 30 giây} (2.400 bước). Checkpoint tốt
nhất (epoch 3) được lưu lại để đánh giá.

\subsection{Phương Pháp 3: Phát Hiện Thống Kê Dựa Trên Perplexity}

Dựa trên giả thuyết rằng LLM gán perplexity thấp hơn cho văn bản do
máy tạo ra (vì nó khớp với phân phối ngôn ngữ của mô hình), chúng tôi
sử dụng \texttt{Qwen/Qwen2.5-1.5B}~\cite{qwen2024qwen25} để tính PPL
của từng bài đăng:

\[
\text{PPL}(x) = \exp\!\left(\frac{1}{T}\sum_{t=1}^{T} \mathcal{L}(x_t \mid x_{<t})\right)
\]

trong đó $\mathcal{L}$ là cross-entropy loss và $T$ là số token. Điểm
phát hiện được định nghĩa là $s = -\text{PPL}(x)$ (điểm cao hơn $\Rightarrow$
khả năng do máy tạo cao hơn). Phương pháp này hoàn toàn zero-shot, không
cần nhãn. Metric đánh giá: AUC-ROC (không phụ thuộc ngưỡng). Cấu hình:
\texttt{bfloat16}, \texttt{device\_map='auto'}, max\_length=512, batch=8.

\subsection{Phương Pháp 4: Học Chuyển Giao (Transfer Learning)}

Để đánh giá khả năng tổng quát hóa liên ngôn ngữ, chúng tôi thực hiện
thí nghiệm \textit{leave-one-language-out}: với mỗi ngôn ngữ, huấn luyện
XLM-RoBERTa trên ba ngôn ngữ còn lại và kiểm tra trên ngôn ngữ bị giữ lại.
Mô hình được khởi tạo mới cho mỗi thí nghiệm. Các siêu tham số chính:

\begin{itemize}
    \item Số epoch: 5; Learning rate: $2 \times 10^{-5}$; Batch size: 16
    \item Max sequence length: 512 tokens; Weight decay: 0,01
    \item Warmup ratio: 0,1; Early stopping patience: 2
    \item Precision: fp16; Train/val split: 80/20
\end{itemize}

Quy trình này mô phỏng kịch bản thực tế khi một ngôn ngữ mới xuất hiện
mà không có dữ liệu huấn luyện, buộc hệ thống phải dựa vào kiến thức
từ các ngôn ngữ khác.
